\sequence{Functions: representations and reference functions}{Functions: representations}{23 November--4 December 2026}
\section{Function, image and preimage}
\begin{cours}[Definition]
A real-valued function associates to each real number $x$ of its domain $D_f$ a unique real number $f(x)$, called the image of $x$. A preimage of $y$ is a real number $x\in D_f$ such that $f(x)=y$. We write $f:x\mapsto f(x)$. A number may have no preimage, one, or several.
\end{cours}
\ex{Let $g(x)=x^2-4x+32$ on $\R$. Calculate $g(-5)$. Let $h(x)=3x+5$; determine the preimage of $6$ under $h$.}
\reponse[3.2]{$g(-5)=25+20+32=77$: $77$ is the image of $-5$, and $-5$ is a preimage of $77$. For $h$, $3x+5=6\iff x=1/3$. The only preimage is $1/3$.}
\ex{For the square function, calculate the images of $-3$ and $3$, then the preimages of $9$, $0$ and $-1$.}
\reponse[3.2]{$f(-3)=f(3)=9$. The preimages of $9$ are $-3,3$; $0$ has a single preimage, $0$; $-1$ has none in $\R$.}
\begin{lstlisting}
def g(x):
    return x*x - 4*x + 32
\end{lstlisting}
\ex{Distinguish \texttt{g(-5)} from \texttt{x = -5}.}
\reponse[2]{The first calls the function and returns $77$. The second assigns $-5$ to the variable \texttt{x}, without calculating the image under $g$.}
\fiche{Domain of definition}
\begin{cours}[Conditions for existence]
$D_f$ contains the values for which the process defining $f$ makes sense. A denominator must be nonzero; a number under a square root must be positive or zero. Also respect any domain imposed by the statement or the context.
\end{cours}
\ex{Determine the domains of definition:
\[f(x)=3x^2-5x+2,\quad g(x)=\frac{2x+1}{x-4},\quad h(x)=\frac{x}{x^2-9},\quad k(x)=\sqrt{2x-6}.\]}
\reponse[5]{$D_f=\R$. $D_g=]-\infty;4[\cup]4;+\infty[$ since $x-4\ne0$. $D_h=]-\infty;-3[\cup]-3;3[\cup]3;+\infty[$ since $x^2-9\ne0$. $D_k=[3;+\infty[$ since $2x-6\geq0$.}
\vigilance{The square root of a negative number does not exist in $\R$; it is not ``a negative square root'' that is excluded.}
\ex[Modelling]{A rectangle of perimeter $20$ m has width $x$ m. Express its area and state the domain if both dimensions must be strictly positive.}
\reponse[5]{The length is $10-x$, so $A(x)=x(10-x)$ in m$^2$. $x>0$ and $10-x>0$ give $D_A=]0;10[$. If degenerate rectangles were allowed, we could take $[0;10]$.}
\fiche{Table of values and representative curve}
\begin{cours}[Definition]
$\mathcal C_f$ is the set of points $(x;f(x))$, with $x\in D_f$. Thus $M(a;b)\in\mathcal C_f\iff a\in D_f$ and $b=f(a)$.
\end{cours}
\ex{On $[-1;4]$, $f(x)=x^2-2x-1$. Complete the table, then draw the curve over the whole interval.}
\begin{center}\renewcommand{\arraystretch}{1.6}\begin{tabular}{c|rrrrrrrr}
$x$&$-1$&$-0.5$&$0$&$1$&$1.5$&$2$&$3$&$4$\\\hline
$f(x)$&\trou{$2$}&\trou{$0.25$}&$-1$&$-2$&\trou{$-1.75$}&$-1$&\trou{$2$}&\trou{$7$}
\end{tabular}\end{center}
\begin{center}\begin{tikzpicture}\begin{axis}[coursaxe,xmin=-2,xmax=5,ymin=-3,ymax=8,width=12cm,height=9cm,xtick={-2,...,5},ytick={-3,...,8}]
\ifprof\addplot[thick,domain=-1:4]{x*x-2*x-1};\addplot[only marks]coordinates{(-1,2)(-.5,.25)(0,-1)(1,-2)(1.5,-1.75)(2,-1)(3,2)(4,7)};\fi
\end{axis}\end{tikzpicture}\end{center}
\reponse[2]{Images: $2$; $0.25$; $-1$; $-2$; $-1.75$; $-1$; $2$; $7$. Plot the points then draw the parabola. Segments joining only these points would give just an approximation.}
\fiche{Graphical solving of equations and inequalities}
\begin{cours}[Method]
The solutions of $f(x)=k$ are the abscissae of the intersections of $\mathcal C_f$ with $y=k$. For $f(x)\leq k$, read the abscissae where the curve is below this line, intersections included; with $<$, the intersections are excluded.
\end{cours}
\ex{The curve represents $f(x)=1+\frac14(x+3)(x+1)(x-2)$ on $[-4;3]$. Solve graphically $f(x)=1$, $f(x)\leq1$ and $f(x)>1$.}
\begin{center}\begin{tikzpicture}\begin{axis}[coursaxe,xmin=-4.2,xmax=3.2,ymin=-4,ymax=7,xtick={-4,...,3},ytick={-4,-2,0,1,3,5,7},width=12cm,height=8cm]
\addplot[very thick,domain=-4:3]{1+(x+3)*(x+1)*(x-2)/4};\addplot[dashed,domain=-4:3]{1};\node[anchor=south east]at(axis cs:3,1){$y=1$};
\end{axis}\end{tikzpicture}\end{center}
\reponse[4]{$f(x)=1$: $-3,-1,2$. $f(x)\leq1$: $[-4;-3]\cup[-1;2]$. $f(x)>1$: $]-3;-1[\cup]2;3]$. The endpoints $-4,3$ come from the domain. The factorisation checks the intersections exactly; the algebraic study of the sign will be consolidated in sequence 8.}
\prof{Graph rebuilt to replace an ambiguous reading in the original booklet. Distinguish graphical conjecture from exact validation.}
\fiche{Comparing two functions; reading a sign}
\begin{cours}[Method]
$f(x)=g(x)$ is read at the intersections of the curves. For $f(x)<g(x)$, $\mathcal C_f$ is strictly below $\mathcal C_g$ at the same abscissa. The sign of $f(x)$ is read relative to the x-axis.
\end{cours}
\ex{On $[-3;3]$, $f(x)=x^2$ and $g(x)=x+2$. Solve $f=g$, $f<g$, $f\geq g$ and determine the sign of $g$.}
\begin{center}\begin{tikzpicture}\begin{axis}[coursaxe,xmin=-3.2,xmax=3.2,ymin=-2,ymax=10,width=12cm,height=8cm,xtick={-3,...,3},ytick={-2,0,2,4,6,8,10}]
\addplot[very thick,domain=-3:3]{x*x};\addplot[thick,dashed,domain=-3:3]{x+2};\node at(axis cs:2.6,8.5){$\mathcal C_f$};\node at(axis cs:2.6,4.2){$\mathcal C_g$};
\end{axis}\end{tikzpicture}\end{center}
\reponse[4]{$f=g$ for $x=-1,2$; $f<g$ on $]-1;2[$; $f\geq g$ on $[-3;-1]\cup[2;3]$. $g$ is strictly negative on $[-3;-2[$, zero at $-2$, strictly positive on $]-2;3]$. Sign and variations are different pieces of information.}
\fiche{Affine functions and absolute value}
\begin{cours}[Definitions]
An affine function is $f(x)=mx+p$ on $\R$. It is linear if $p=0$, constant if $m=0$. For $m\ne0$, it vanishes at $-p/m$; its sign is that of $m$ to the right of this number.
The function $x\mapsto|x|$ is defined on $\R$, positive or zero, and vanishes only at $0$.
\end{cours}
\ex{Classify (several answers possible) as affine, linear, constant: $5x+2$; $5x$; $7+2x-7$; $6$; $6(4x-2)$; $6x+5-6x$.}
\reponse[2.4]{All are affine. $5x$ and $2x$ are linear; $6$ and $5$ are constant; $24x-12$ is neither linear nor constant.}
\begin{center}\begin{tikzpicture}\begin{axis}[coursaxe,xmin=-4,xmax=4,ymin=-1,ymax=4,width=9cm,height=5cm,xtick={-4,...,4},ytick={0,...,4}]
\addplot[thick,domain=-4:4]{abs(x)};\node at(axis cs:2.7,3.7){$y=|x|$};\end{axis}\end{tikzpicture}\end{center}
\begin{cours}[Equations and inequalities]
If $k<0$, $|x|=k$ and $|x|\leq k$ have no solution, and $|x|\geq k$ is true on $\R$. If $k=0$, the first two give $x=0$, the third is true on $\R$. For $k>0$:
\[|x|=k\iff x\in\{-k;k\},\quad |x|\leq k\iff x\in[-k;k],\]
\[|x|\geq k\iff x\in]-\infty;-k]\cup[k;+\infty[.\]
\end{cours}
\ex{Solve $|x|=5$, $|x|=-3$, $|x|\leq3$, $|x|\geq4$.}
\reponse[1.6]{$\{-5;5\}$; $\varnothing$; $[-3;3]$; $]-\infty;-4]\cup[4;+\infty[$.}
\fiche{Square and inverse functions}
\begin{cours}[Definitions and sign]
$x\mapsto x^2$ is defined on $\R$; its curve is a parabola. $x^2\geq0$, with equality only at $0$.
$x\mapsto1/x$ is defined on $\R^*$; its curve is a hyperbola. $1/x$ has the sign of $x$ and never vanishes.
\end{cours}
\begin{center}\begin{tikzpicture}\begin{axis}[coursaxe,xmin=-3,xmax=3,ymin=-1,ymax=9,width=8cm,height=5cm,xtick={-2,0,2}]
\addplot[thick,domain=-3:3]{x*x};\node at(axis cs:1,6){$y=x^2$};\end{axis}\end{tikzpicture}
\begin{tikzpicture}\begin{axis}[coursaxe,xmin=-4,xmax=4,ymin=-4,ymax=4,width=8cm,height=5cm,xtick={-4,-2,0,2,4}]
\addplot[thick,domain=-4:-.25]{1/x};\addplot[thick,domain=.25:4]{1/x};\node at(axis cs:2.5,2){$y=1/x$};\end{axis}\end{tikzpicture}\end{center}
\begin{cours}[Equations and inequalities with $x^2$]
If $k<0$, $x^2=k$ and $x^2\leq k$ have no solution; $x^2\geq k$ is true on $\R$. If $k=0$, the first two give $x=0$, the third is true on $\R$. For $k>0$:
\[x^2=k\iff x\in\{-\sqrt k;\sqrt k\},\quad x^2\leq k\iff x\in[-\sqrt k;\sqrt k],\]
\[x^2\geq k\iff x\in]-\infty;-\sqrt k]\cup[\sqrt k;+\infty[.\]
\end{cours}
\ex{Solve $x^2=4$, $x^2=-1$, $x^2\geq4$, $x^2\geq-1$, $1/x=4$, $1/x=0$.}
\reponse[3.2]{$\{-2;2\}$; $\varnothing$; $]-\infty;-2]\cup[2;+\infty[$; $\R$; $\{1/4\}$; $\varnothing$. More generally, $1/x=k$ has the unique solution $1/k$ if $k\ne0$, none if $k=0$.}
\fiche{Inequalities with the inverse function}
\begin{methode}
Respect $x\ne0$ and read both branches of the hyperbola. Do not multiply an inequality by $x$ without knowing its sign.
\end{methode}
\begin{center}\small\renewcommand{\arraystretch}{2.5}\begin{tabular}{c|c|c|c}
&$k<0$&$k=0$&$k>0$\\\hline
$1/x\leq k$&$[1/k;0[$&$]-\infty;0[$&$]-\infty;0[\cup[1/k;+\infty[$\\
$1/x\geq k$&$]-\infty;1/k]\cup]0;+\infty[$&$]0;+\infty[$&$]0;1/k]$
\end{tabular}\end{center}
\ex{Solve $1/x\geq2$, $1/x\leq-1$, $1/x<1/2$ and $1/x>-2$.}
\reponse[5]{$]0;1/2]$; $[-1;0[$; $]-\infty;0[\cup]2;+\infty[$; $]-\infty;-1/2[\cup]0;+\infty[$. A strict inequality excludes the intersection with $y=k$; $0$ is always forbidden.}
\ex{A student writes: ``$1/x\leq2\iff1\leq2x$, so $x\geq1/2$''. What is missing?}
\reponse[4]{They implicitly assumed $x>0$. For $x<0$, $1/x<0<2$: these values also work. The correct set is $]-\infty;0[\cup[1/2;+\infty[$. The number $-1$ refutes their equivalence.}
\fiche{Using square root and cube}
\begin{cours}[Definitions]
$x\mapsto\sqrt x$ is defined on $[0;+\infty[$ and takes values that are positive or zero. For $k\geq0$, $\sqrt x=k\iff x=k^2$; for $k<0$, no solution.
$x\mapsto x^3$ is defined on $\R$. Every real number $k$ has a unique preimage, written $\sqrt[3]{k}$.
\end{cours}
\begin{center}\begin{tikzpicture}\begin{axis}[coursaxe,xmin=-1,xmax=9,ymin=-1,ymax=4,width=8cm,height=5cm,xtick={0,1,4,9}]
\addplot[thick,domain=0:9]{sqrt(x)};\node at(axis cs:5,3){$y=\sqrt x$};\end{axis}\end{tikzpicture}
\begin{tikzpicture}\begin{axis}[coursaxe,xmin=-2,xmax=2,ymin=-8,ymax=8,width=8cm,height=5cm,xtick={-2,-1,0,1,2}]
\addplot[thick,domain=-2:2]{x*x*x};\node at(axis cs:-1,5){$y=x^3$};\end{axis}\end{tikzpicture}\end{center}
\ex{Solve $\sqrt x=3$; $\sqrt x=-2$; $\sqrt x\leq2$; $\sqrt x>3$; $x^3=-8$; $x^3\leq27$.}
\reponse[4]{$\{9\}$; $\varnothing$; $[0;4]$; $]9;+\infty[$; $\{-2\}$; $]-\infty;3]$. Respect the domain for square roots. The cube function preserves order.}
\ex{Compare $x^2$ and $x$ for $x=1/2$, then $x=2$. Is the square of a positive number always greater than the number itself?}
\reponse[2.4]{$1/4<1/2$ but $4>2$. The claim is false; $1/2$ is a counterexample. General study in sequence 10.}
